\cotefile{59}
\scanpage{2}\mspage{1}

\diagram{\begin{tikzcd}[ampersand replacement=\&,row sep=small]
X\arrow[d,"f"']\\ S
\end{tikzcd}}

$f$ propre plat présentation finie \quad $f_*(\underline O_X)\simeq\underline O_S$ univ!

\diagram{\begin{tikzcd}[ampersand replacement=\&]
X\arrow[r,dash]\arrow[d,dash]\&X_P\arrow[d,dash]\\
S\arrow[r,dash]\&P
\end{tikzcd}}

\[
\underline{\mathrm{Pic}}_{X/S}=P
\]

a) Supposons $X/S$ muni d'une section $g$. Alors sur $X\times_S P$ il y a un faisceau inversible canonique $\mathcal L_g$ [trivialisé sur $g(S)\times_S P$]. On en conclut

\[
(Rf_P)_*(\mathcal L_g)\in D_{\mathrm{parf}}(P)
\qquad [f_P:X\times_S P\mathbin{-}P].
\]

Nous nous intéressons en particulier à

% slanted
\margin{on suppose \unsure{th. de finitude|th. du préambule} \unsure{valables}, p. ex. $f$ projectif ou $S$ loc. noeth. etc.}

\[
\det{}^{*}(Rf_P)_*(\mathcal L_g)=M_g
\]

qui est un module inversible [gradué s. il y a lieu] sur $P$.

b) Remplaçons $g$ par une section $g'$. On a alors un isom canonique

\[
\mathcal L_{g'}=\mathcal L_g\otimes_{\underline O_P}{g'}_P^{*}(\mathcal L_g)^{-1}
\]

d'où

\[
(Rf_P)_*(\mathcal L_{g'})\simeq(Rf_P)_*(\mathcal L_g)\,{g'}_P^{*}(\mathcal L_g)^{-1}
\]

et par suite un isom canonique

\[
M_{g'}\simeq M_g\,{g'}_P^{*}(\mathcal L_g)^{-\chi}=M_g\left(N_gN_{g'}^{-1}\right)^{\chi}.\tag*{(*)}
\]

où $\chi$ est la fonction continue $P\longrightarrow\mathbb Z$ donnée par le rang de $Rf_{P*}(\mathcal L_g)$, i.e. la car. d'EP des faisceaux inversibles sur les fibres de $X$ \ldots.., et où $N$ désigne un faisceau inv. \inserted{``de Weil''} quelconque sur $X\times_{\unsure{P|S}}P$, et $N_g$ (resp. $N_{g'}$) les faisceaux inv. sur $P$ qu'il définit via la section $g$ (resp. $g'$). On notera que le facteur de droite

% margin note written sideways
\margin{$N_gN_{g'}^{-1}$ est \uline{algébriquement équivalent} à zero sur $P$ rel/$S$, donc $M_{g'}$ est alg. équiv. à $M_g$ \unsure{rel/$S$|polaris\ldots}.}

c) Soit $\delta$ une section de $\underline{\mathrm{NS}}_{X/S}$, qui définit donc un torseur $P^\delta$ sous $P^{\unsure{\circ|0}}=\underline{\mathrm{Pic}}^\circ_{X/S}$, et considérons la restriction de $M_g$, \inserted{$M_{g'},\ldots$} à $P^\delta$. Nous supposons pour simplifier, par la suite, que $P^\circ$ est un schéma \uline{abélien} sur $S$.
